% Part: lambda-calculus
% Chapter: introduction
% Section: beta

\documentclass[../../../include/open-logic-section]{subfiles}

\begin{document}

\olfileid{lam}{int}{bet}
\olsection{\texorpdfstring{$\beta$-归约}{beta-归约}}

当我们看到 $(\lambd[m][(\lambd[y][y]) m])$ 时，自然会猜想它与 $\lambd[m][m]$ 有某种联系，即第二个项应为第一个项“简化”的结果。\emph{$\beta$-归约}的概念形式化了这一直觉。

\begin{defn}[$\beta$-收缩, $\bredone$] \ollabel{defn:betacontr}
  \emph{$\beta$-收缩}（$\bredone$）是项上满足以下条件的最小相容关系：
  \[
    (\lambd[x][N])Q \bredone \Subst{N}{Q}{x} 
  \]
  若 $P \bredone Q$，则称 $P$ \emph{$\beta$-收缩}到 $Q$。
  形如 $(\lambd[x][N])Q$ 的项称为\emph{可约式}。
\end{defn}

\begin{prob} \ollabel{prob:def}
  如同我们在 \olref[lam][syn][alp]{defn:aconvone} 中对约束变元更改所做的那样，写出 $\beta$-收缩的等价归纳定义。
\end{prob}
  
\begin{defn}[$\beta$-归约, $\bred$] \ollabel{defn:betared}
  \emph{$\beta$-归约}（$\bred$）是包含 $\bredone$ 的最小自反、传递关系。
  若 $P \bred Q$，则称 $P$ $\beta$-归约到 $Q$。
\end{defn}

在上下文明确时，我们将用 $\redone$ 代替 $\bredone$，用 $\red$ 代替 $\bred$。

非形式地说，$M \bred N$ 当且仅当 $M$ 可以通过零步或多步 $\beta$-收缩变为 $N$。

\begin{defn}[$\beta$-范式]
一个不能再进行 $\beta$-收缩的项称为\emph{$\beta$-范式}。
\end{defn}

如果 $M \bred N$ 且 $N$ 是 $\beta$-范式，则称 $N$ 是~$M$ 的一个\emph{范式}。人们可能会问一个项的范式是否唯一，答案是肯定的，后文将看到这一点。

我们来看几个例子。
\begin{enumerate}
\item 我们有
\begin{align*}
(\lambd[x][xxy]) \lambd[z][z] & \redone (\lambd[z][z])(\lambd[z][z]) y \\
& \redone (\lambd[z][z]) y \\
& \redone y
\end{align*}
\item “简化”一个项实际上可能使其更复杂：
\begin{align*}
(\lambd[x][xxy])(\lambd[x][xxy]) & \redone (\lambd[x][xxy])(\lambd[x][xxy])y \\
& \redone (\lambd[x][xxy])(\lambd[x][xxy])yy \\
& \redone \dots
\end{align*}
\item 它也可能使项保持不变：
\[
(\lambd[x][xx])(\lambd[x][xx]) \redone (\lambd[x][xx])(\lambd[x][xx])
\]
\item 此外，有些项可以以不止一种方式归约；例如，
\[
(\lambd[x][(\lambd[y][yx]) z]) v \redone (\lambd[y][yv]) z
\]
通过收缩最外层的应用；以及
\[
(\lambd[x][(\lambd[y][yx]) z]) v \redone (\lambd[x][zx]) v
\]
通过收缩最内层的应用。不过注意，在这种情况下，两个项都进一步归约到了同一个项 $zv$。
\end{enumerate}

最后一个例子中的最终结果并非巧合，而是展示了λ演算的一个深刻而重要的性质，即 邱奇--罗瑟 性质（丘奇-罗瑟性质）。

\begin{digress}
  一般来说，$\beta$-归约一个项的方式不止一种，因此人们发明了许多归约策略，其中最常见的是\emph{自然策略}。自然策略总是收缩\emph{最左}可约式，其中可约式的位置定义为它在项中的起始点。自然策略有一个有用的性质：一个项能通过某种策略归约为范式，当且仅当它能使用自然策略归约为范式。在下文中，除非另有说明，我们将使用自然策略。
\end{digress}

\begin{defn}[$\beta$-等价, $\equal$]
  \emph{$\beta$-等价}（$\equal$）是通过以下规则归纳定义的关系：
  \begin{enumerate}
  \item $M \equal M$。
  \item 如果 $M \equal N$，则 $N \equal M$。
  \item 如果 $M \equal N$，$N \equal O$，则 $M \equal O$。
  \item 如果 $M \equal N$，则 $PM \equal PN$。
  \item 如果 $M \equal N$，则 $MQ \equal NQ$。
  \item 如果 $M \equal N$，则 $\lambd[x][M] \equal \lambd[x][N]$。
  \item $(\lambd[x][N])Q \equal \Subst{N}{Q}{x}$。
  \end{enumerate}
\end{defn}

前三条规则使该关系成为等价关系；接下来的三条使其兼容；最后一条确保它包含 $\beta$-收缩。

非形式地说，两个项是 $\beta$-等价的，当且仅当其中一个可以通过零次或多次 $\beta$-收缩或 $\beta$-收缩的逆变为另一个。$\beta$-收缩的逆定义如下：$M$ 逆$\beta$-收缩到 $N$ 当且仅当 $N$ $\beta$-收缩到 $M$。

除了上述规则，我们还将用更多规则扩展这一关系，并用 $\equal[X]$ 表示扩展后的等价关系，其中 $X$ 是扩展规则。

\end{document}
